\section*{Conclusion}

Ce travail a permis d'établir un modèle homogénéisé décrivant le
transport de matière et de charge au sein d'une électrode poreuse.
À partir des équations locales de conservation, de Nernst--Planck,
de Poisson et d'Ohm, une procédure d'homogénéisation spatiale a été
mise en œuvre afin de remplacer la géométrie complexe du milieu poreux
par un milieu homogène équivalent caractérisé par des propriétés
effectives.

L'analyse des différentes hypothèses simplificatrices a ensuite montré
comment ce modèle général pouvait être réduit à des formulations plus
simples selon les phénomènes retenus. En particulier, le modèle
classique fondé sur le module de Thiele apparaît comme un cas limite
correspondant à un transport purement diffusif couplé à une réaction
du premier ordre.

Le cadre obtenu reste cependant plus général que l'approche de Thiele,
car il permet de conserver les effets de migration électrique,
de convection, de transport électronique et de structure poreuse.
Il constitue ainsi une base de travail pour l'étude de situations où
les hypothèses habituellement utilisées deviennent insuffisantes.

Plusieurs perspectives peuvent être envisagées pour la suite de ce
travail. La première consiste à résoudre numériquement le système
homogénéisé complet afin d'évaluer l'influence relative des différents
mécanismes de transport. Une autre perspective concerne l'amélioration
de la description des réactions électrochimiques, notamment en
remettant en question certaines hypothèses cinétiques utilisées dans
ce document et en introduisant des modèles plus adaptés aux conditions
de fonctionnement considérées.